\chapter{Introdução}

O relatório final deverá ter um máximo de 80 páginas (em relação aos capítulos numerados, sem considerar bibliografia, anexos, listas, \dots).

Fornecem-se aqui algumas orientações genéricas sobre a organização do conteúdo. Em caso de detectarem alguma informação contraditória, prevalece sempre o que estiver indicado
no site da Faculdade de Ciências, na secção sobre \href{https://ciencias.ulisboa.pt/pt/admissao-provas-academicas-2-ciclo}{Admissão a Provas Académicas de 2º ciclo}.

O seu conteúdo deve realçar o trabalho realizado pelo aluno e a sua contribuição concreta no trabalho. Por exemplo, se o trabalho consiste no desenvolvimento de vários módulos a serem integrados num sistema mais global, o aluno deverá preocupar-se em descrever a parte que desenvolveu, como desenvolveu, que ferramentas usou, que alternativas poderiam existir, etc., em vez de efectuar uma descrição exaustiva das funcionalidades de todo o sistema.

O número de capítulos no relatório final não é rígido. No entanto, recomenda-se que sejam adoptados os seguintes princípios para a organização do relatório:

\begin{enumerate}
\item Um capítulo \emph{introdutório} no qual se apresentam o contexto do trabalho, se resume o trabalho desenvolvido, se identificam as contribuições deste e se apresenta a estrutura do próprio relatório. Deverá também ser mencionado sucintamente o enquadramento institucional em que o trabalho decorreu.
\item Um capítulo no qual se apresentam \emph{em pormenor} os \emph{objectivos} do trabalho, o \emph{contexto subjacente}, a \emph{metodologia} utilizada no seu desenvolvimento bem como o \emph{planeamento} efectuado para o concretizar. Deve também ser apresentada uma confrontação com o plano de trabalho inicial analisando as razões de eventuais desvios ocorridos.
\item Um capítulo onde é descrito o \emph{trabalho realizado}. Este é um dos capítulos fundamentais do relatório. Apresenta concretamente o que se fez de facto e as ferramentas usadas. De notar que é importante que fique claro qual a contribuição concreta do trabalho, sobretudo em casos de trabalho em equipa. Neste capítulo poderão ser inseridas questões relevantes da área de estudo em que o trabalho se integra, assim como o possível enquadramento num trabalho mais amplo. Eventualmente, e em função do âmbito e dimensão do trabalho, este capítulo poderá ser substituído por um conjunto de outros capítulos, que englobem em si o \emph{trabalho relacionado}, a \emph{análise} do problema, o \emph{desenho} da solução, a \emph{implementação} da solução e a \emph{avaliação} desta.
\item Um capítulo no qual são apresentadas as \emph{conclusões}. Para além de um \emph{sumário} do trabalho realizado, deve ser feito um \emph{comentário crítico} e serem apresentadas possibilidades de \emph{trabalho futuro} referindo o que falta fazer e o que poderá ser melhorado.
\item Um capítulo com a \emph{bibliografia} - lista de documentos usados e outras referências consideradas relevantes.
\item Um conjunto de capítulos com os \emph{anexos}. Quaisquer listagens, informação confidencial ou outras descrições muito pormenorizadas não devem ser integradas no corpo principal do relatório. Se houver necessidade de as apresentar, sugere-se a sua introdução em anexos ou em documentos separados. Os anexos suplementam o relatório e como tal devem ser referidos no corpo principal do relatório, descrevendo o tipo de informação que se detalha em anexo.
\item Poderá incluir uma lista de Acrónimos (como \acrshort{cpu} ou \acrshort{mmu}) e um Glossário (para descrever, por exemplo, \emph{\gls{cmd}}).
\end{enumerate}

Quando concluir a sua Dissertação, ou Relatório Final, deverá proceder de acordo com o indicado
no \href{https://ciencias.ulisboa.pt/pt/guias-estudantes}{Guia Académico} do ano lectivo em causa ou a informação disponilizada no site da FCUL, na secção \href{https://ciencias.ulisboa.pt/pt/admissao-provas-academicas-2-ciclo}{Admissão a Provas Académicas de 2º ciclo}.

\section{Motivação}

Lorem \index{Lorem} ipsum dolor sit amet, consectetuer adipiscing elit, sed diam nonummy nibh euismod tincidunt ut laoreet dolore magna aliquam erat volutpat. Ut wisi enim ad minim veniam, quis nostrud exerci tation ullamcorper suscipit lobortis nisl ut aliquip ex ea commodo consequat. Duis autem vel eum iriure dolor in hendrerit in vulputate velit esse molestie consequat, vel illum dolore eu feugiat nulla facilisis at vero eros et accumsan et iusto odio dignissim qui blandit praesent luptatum zzril delenit augue duis dolore te feugait nulla facilisi. Nam liber tempor cum soluta nobis eleifend option congue nihil imperdiet doming id quod mazim placerat facer possim assum. Typi non habent claritatem insitam; est usus legentis in iis qui facit eorum claritatem. Investigationes demonstraverunt lectores legere me lius quod ii legunt saepius. Claritas est etiam processus dynamicus, qui sequitur mutationem consuetudium lectorum. Mirum est notare quam littera gothica, quam nunc putamus parum claram, anteposuerit litterarum formas humanitatis per seacula quarta decima et quinta decima. Eodem modo typi, qui nunc nobis videntur parum clari, fiant sollemnes in futurum. Lorem ipsum dolor sit amet, consectetuer adipiscing elit, sed diam nonummy nibh euismod tincidunt ut laoreet dolore magna aliquam erat volutpat. Ut wisi enim ad minim veniam, quis nostrud exerci tation ullamcorper suscipit lobortis nisl ut aliquip ex ea commodo consequat. Duis autem vel eum iriure dolor in hendrerit in vulputate velit esse molestie consequat, vel illum dolore eu feugiat nulla facilisis at vero eros et accumsan et iusto odio dignissim qui blandit praesent luptatum zzril delenit augue duis dolore te feugait nulla facilisi. Nam liber tempor cum soluta nobis eleifend option congue nihil imperdiet doming id quod mazim placerat facer possim assum. Typi non habent claritatem insitam; est usus legentis in iis qui facit eorum claritatem. Investigationes demonstraverunt lectores legere me lius quod ii legunt saepius.

\section{Objectivos}

\blindtext[1]

\blindtext[2]

\section{Contribuições}

\Blindtext[2]

\section{Estrutura do documento}

Este documento está organizado da seguinte forma:
\begin{itemize}
\item Capítulo 2 – AAA
\item Capítulo 3 – BBB
\end{itemize}

